\documentclass[10pt]{article}
\providecommand{\FIGDIR}{fig}
\newcommand{\DIGESTDATE}{22 August 2026}
\newcommand{\PERIODSTART}{16 August 2026}
\newcommand{\PERIODEND}{22 August 2026}
% ===================== PM Research Watch - shared preamble =====================
\usepackage[T1]{fontenc}
\usepackage[utf8]{inputenc}
\usepackage{amsmath,amssymb}
\usepackage{tgpagella}
\usepackage{tgheros}
\renewcommand{\sfdefault}{qhv}
\renewcommand{\ttdefault}{lmtt}
\usepackage{microtype}
\usepackage[english]{babel}

\usepackage[a4paper,top=17mm,bottom=17mm,left=17mm,right=17mm,
            headsep=6mm,footskip=10mm]{geometry}
\usepackage{graphicx}
\usepackage{xcolor}
\usepackage{colortbl}
\usepackage{tikz}
\usetikzlibrary{calc}
\usepackage[most]{tcolorbox}
\usepackage{titlesec}
\usepackage{fancyhdr}
\usepackage{booktabs}
\usepackage{array}
\usepackage{longtable}
\usepackage{enumitem}
\usepackage{multicol}
\usepackage{ragged2e}
\usepackage[hidelinks]{hyperref}

% ------------------------------------------------------------------ palette
\definecolor{Ink}    {HTML}{14262E}
\definecolor{Deep}   {HTML}{0F4C5C}
\definecolor{Teal}   {HTML}{1F7A8C}
\definecolor{Sky}    {HTML}{4FA3C4}
\definecolor{Sage}   {HTML}{6FA287}
\definecolor{Amber}  {HTML}{E3A008}
\definecolor{Clay}   {HTML}{B5643C}
\definecolor{Coral}  {HTML}{D1495B}
\definecolor{Violet} {HTML}{6A5B9E}
\definecolor{Slate}  {HTML}{4A6572}
\definecolor{Paper}  {HTML}{FBF9F5}
\definecolor{Mist}   {HTML}{EEEAE2}
\definecolor{Rule}   {HTML}{D9D3C9}

\pagecolor{Paper}
\color{Ink}

\hypersetup{colorlinks=true, linkcolor=Deep, urlcolor=Teal, citecolor=Teal,
            pdftitle={PM Research Watch}, pdfauthor={Automated daily digest}}

% ------------------------------------------------------------------ headings
\titleformat{\section}
  {\sffamily\bfseries\large\color{Deep}}{}{0pt}
  {}[{\vspace{1.2mm}\color{Amber}\titlerule[1.1pt]}]
\titlespacing*{\section}{0pt}{5.5mm}{2.8mm}

\titleformat{\subsection}
  {\sffamily\bfseries\normalsize\color{Teal}}{}{0pt}{}
\titlespacing*{\subsection}{0pt}{3.6mm}{1.4mm}

\setlist[itemize]{leftmargin=4.2mm, itemsep=1.1mm, topsep=1.3mm, parsep=0pt,
                  label=\textcolor{Teal}{\footnotesize$\blacktriangleright$}}

% ------------------------------------------------------------------ header/footer
\pagestyle{fancy}
\fancyhf{}
\renewcommand{\headrulewidth}{0pt}
\renewcommand{\footrulewidth}{0pt}
\fancyhead[L]{\sffamily\scriptsize\color{Slate}PM RESEARCH WATCH\;\textcolor{Rule}{\textbar}\;\DIGESTDATE}
\fancyhead[R]{\sffamily\scriptsize\color{Slate}Particulate matter monitoring \& health effects}
\fancyfoot[C]{\sffamily\scriptsize\color{Slate}\thepage}

% ------------------------------------------------------------------ components
% metric tile
\newcommand{\metric}[3]{%
  \begin{tcolorbox}[enhanced, colback=white, colframe=Rule, boxrule=0.5pt,
    arc=1.6mm, left=2.2mm, right=2.2mm, top=1.8mm, bottom=1.8mm,
    borderline west={1.7mm}{0pt}{#3}, nobeforeafter, valign=center,
    height=17mm, width=\linewidth]
    \RaggedRight
    {\sffamily\bfseries\fontsize{17}{17}\selectfont\color{#3}#1}\\[0.5mm]
    {\sffamily\fontsize{6.5}{7.4}\selectfont\color{Slate}#2}
  \end{tcolorbox}}

% pull-quote / key insight box
\newtcolorbox{keybox}[1][]{enhanced, colback=Mist, colframe=Mist,
  arc=1.6mm, left=3.2mm, right=3.2mm, top=2.4mm, bottom=2.4mm,
  borderline west={1.7mm}{0pt}{Amber}, #1}

% subtopic band
\newcommand{\band}[2]{%
  \vspace{3.2mm}\noindent
  \begin{tcolorbox}[enhanced, colback=#2, colframe=#2, arc=1.2mm,
    left=2.6mm, right=2.6mm, top=1.1mm, bottom=1.1mm, boxrule=0pt,
    width=\linewidth, nobeforeafter]
    {\sffamily\bfseries\small\color{white}#1}
  \end{tcolorbox}\vspace{1.6mm}}

% one paper entry: {short cite}{title}{journal}{doi}{body}
\newcommand{\paperentry}[5]{%
  \par\vspace{1.9mm}\noindent
  {\sffamily\bfseries\footnotesize\color{Ink}#1}\;%
  {\sffamily\fontsize{7}{8}\selectfont\color{Slate}#3}\\[0.4mm]
  {\footnotesize\itshape\color{Deep}#2}\\[0.6mm]
  {\footnotesize #4}\\[0.3mm]
  {\fontsize{6.6}{7.6}\selectfont\color{Slate}\texttt{doi:}\,\href{https://doi.org/#5}{\texttt{\detokenize{#5}}}}
  \par}

% caption style
\newcommand{\figcap}[1]{\par\vspace{1.0mm}
  {\fontsize{7.2}{8.6}\selectfont\color{Slate}\textsf{\textbf{Fig.}}\; #1\par}\vspace{1.5mm}}

\providecommand{\permil}{\textperthousand}

\setlength{\parindent}{0pt}
\setlength{\parskip}{0pt}
\renewcommand{\arraystretch}{1.18}


\begin{document}

% ==================================================================== MASTHEAD
\thispagestyle{empty}
\begin{tikzpicture}[remember picture, overlay]
  \fill[Deep] (current page.north west) rectangle
      ($(current page.north east)+(0,-67mm)$);
  \fill[Teal] ($(current page.north west)+(0,-67mm)$) rectangle
      ($(current page.north east)+(0,-68.4mm)$);
  \foreach \x/\y/\r in {22/12/0.9, 41/26/1.5, 63/9/0.7, 88/31/1.1, 112/17/1.9,
                        138/29/0.8, 159/13/1.3, 178/54/1.0, 196/21/1.6,
                        30/44/0.6, 74/57/1.2, 125/48/0.9, 168/6/0.8, 191/59/0.7,
                        52/18/1.0, 100/60/0.6, 148/20/1.1, 66/49/1.4, 205/45/0.8}
    \fill[Sky, opacity=0.30] ($(current page.north west)+(\x mm,-\y mm)$) circle (\r mm);
\end{tikzpicture}

\vspace*{4mm}
{\sffamily\fontsize{7.6}{9}\selectfont\color{Sky}%
 WEEKLY ROLLUP\;\textbullet\;PhD RESEARCH WATCH\par}
\vspace{2.2mm}
{\sffamily\bfseries\fontsize{27}{29}\selectfont\color{white}%
 Particulate Matter\par}
{\sffamily\fontsize{27}{29}\selectfont\color{Sky}%
 Week in Review\par}
\vspace{5.0mm}
{\sffamily\fontsize{8.4}{10}\selectfont\color{white}%
 \mbox{Week ending \DIGESTDATE}\;\;\textcolor{Teal}{\textbar}\;\;\mbox{\PERIODSTART{} to \PERIODEND}\;\;
 \textcolor{Teal}{\textbar}\;\;\mbox{112 records pooled from 7 daily issues}\par}
\vspace{18mm}

% -------------------------------------------------------------- metric strip
\noindent
\begin{minipage}[t]{0.190\linewidth}\metric{112}{RECORDS POOLED\\(112 UNIQUE DOIs)}{Deep}\end{minipage}\hfill
\begin{minipage}[t]{0.190\linewidth}\metric{52\%}{ON THE MEASUREMENT\\SIDE (58 RECORDS)}{Teal}\end{minipage}\hfill
\begin{minipage}[t]{0.190\linewidth}\metric{10}{SUBTOPICS\\ACTIVE}{Sky}\end{minipage}\hfill
\begin{minipage}[t]{0.190\linewidth}\metric{41}{POOLED EFFECT\\ESTIMATES}{Amber}\end{minipage}\hfill
\begin{minipage}[t]{0.190\linewidth}\metric{22}{ABSTRACT-FREE\\(20\% OF WEEK)}{Clay}\end{minipage}

\vspace{6mm}

% -------------------------------------------------------------- week in one box
\section*{\sffamily\bfseries\color{Deep}The week in one paragraph}
\vspace{-1.5mm}{\color{Teal}\rule{\linewidth}{1.1pt}}\vspace{2.5mm}

\begin{keybox}
{\sffamily\bfseries\small\color{Deep}Six independent groups stopped analysing PM$_{2.5}$
as a mass this week and started analysing it as a list of constituents --- and in the two
endpoints where a single constituent led, it was the same one. Meanwhile a fifth of the
week's corpus arrived with no abstract, which is the price this pipeline is now paying for
its own coverage.}\\[1.4mm]
\footnotesize
The methodological through-line is constituent decomposition, and it is unusually
consistent. \textbf{Wu et al.} (16 Aug) put prenatal PM$_{2.5}$ and its components through
a mixture model for childhood overweight; \textbf{Wang et al.} (16 Aug) do the same for
early-life anaemia across China, India and 26 African countries; \textbf{Zhao et al.}
(17 Aug) report spontaneous abortion by constituent, with \textbf{ammonium highest
(OR 1.19) --- above total mass (1.18)}; \textbf{Parmar et al.} (18 Aug) separate black
carbon from bulk PM$_{2.5}$ for neuropsychiatric hospitalisation in Delhi and find the
component term nearly an order of magnitude larger; \textbf{Chen et al.} (21 Aug) run
five constituents against incident depression in 24{,}940 UK Biobank participants and
again land on \emph{ammonium}; and \textbf{Wei et al.} (21 Aug) run all fifteen components
against mortality in 528{,}127 lung-cancer patients by weighted quantile sum.

\textbf{The answers do not agree, and that is the finding.} Ammonium and nitrate --- the
secondary inorganic fraction --- lead in the reproductive and psychiatric endpoints, while
Wei's mortality weights fall on silicon, vanadium, zinc and iron, a crustal, traffic and
heavy-fuel-oil signature. Two different exposures are being called PM$_{2.5}$. The week
also supplies the atmospheric-chemistry half of that argument: \textbf{Chai et al.}
(21 Aug) measure secondary inorganic aerosol at \textbf{49.55\% of PM$_{2.5}$ mass} in
suburban Nanjing and show the system is total-nitrate-limited, so the very fraction that
the epidemiology is now singling out is the one that does \emph{not} respond to ammonia
control. If a component drives the risk, mass-only monitoring cannot see it and
mass-targeted policy cannot move it.

\textbf{Against that, an operational fact that should be read as a cost.} 22 of 112
records (20\%) carry no abstract from any checked source and are classified but unread ---
almost double last week's 15 of 134 (11\%). \textit{Metadata only} is now the single
largest architecture label in the pooled corpus, ahead of every real study design.
This follows directly from widening the Crossref by-ISSN sweep: the leg buys coverage
of journals PubMed indexes late or not at all, and pays for it in publishers that
deposit a title and a DOI and nothing more.
\end{keybox}

\clearpage

% -------------------------------------------------------------- trend
\section*{\sffamily\bfseries\color{Deep}Subtopic trend and inflection}
\vspace{-1.5mm}{\color{Teal}\rule{\linewidth}{1.1pt}}\vspace{2.5mm}

\noindent\includegraphics[width=\linewidth]{\FIGDIR/w1_subtopic_trend.png}
\figcap{Per-issue subtopic counts across the week from \texttt{state/metrics.csv}.
\textbf{Read the denominators first}: daily totals were 10, 12, 18, 22, 18, 19 and 13 ---
a range of 2.2$\times$, and a flatter week than 9--15 August (12 to 28). The 16 August
floor of 10 is a Sunday and the smallest issue the watch has recorded; the 22 August
figure of 13 is a Saturday. \textbf{No day in this week carries a publisher batch
artefact of the kind that distorted 11--12 August}, so unlike last week the trend line can
be read more or less at face value.}

\footnotesize
\textbf{The measurement side took the week back.} \textit{Exposure assessment \&
modelling} rose \textbf{26 $\rightarrow$ 36} while \textit{Sensing, forecasting \&
instrumentation} fell \textbf{34 $\rightarrow$ 22}; together they are 58 of 112 records,
a measurement-side share of \textbf{52\% against last week's 45\% and 54\% the week
before}. \textbf{Do not read this as a recovery in instrumentation.} The sensing cluster
shrank; what grew is exposure and deposition work, and four of the eight
\textit{Exposure} records in the 21 August issue alone measure PM by what it lands on ---
pine needles, privet leaves, tree canopies --- rather than by what it does in air. The
oscillation between 54\%, 45\% and 52\% across three weeks is within the range this
pipeline's own harvesting legs can produce, and no claim about the field should be built
on it.

\vspace{1.2mm}
\textbf{Contracting, and worth watching.} \textit{Occupational \& indoor} fell
\textbf{9 $\rightarrow$ 4}, \textit{Mechanistic toxicology} \textbf{11 $\rightarrow$ 8}
and \textit{Cardiovascular \& metabolic} \textbf{7 $\rightarrow$ 3}. The mechanistic drop
breaks a two-week run at exactly 11 and is the first movement in that cluster since the
watch began keeping the count. Cardiovascular at 3 is the lowest recorded for what is
normally the field's largest health cluster --- against 10 respiratory-endpoint records
this week. One week is not a trend; two would be.

\vspace{1.2mm}
\textbf{Holding.} \textit{Other clinical endpoints} rose \textbf{11 $\rightarrow$ 13},
extending last week's finding that the clinical tail is now longer than the
cardiorespiratory trunk, and the content stays unfamiliar: cirrhotic hepatic
decompensation, diabetic microvascular complications, systemic lupus, and --- new this
week --- \textbf{the prognosis of an emerging tick-borne viral infection}. All ten
subtopics were active on at least one day.

%% LAYOUT RULE (2026-08-02): architecture donut in the whitespace here, not on its own page.
\vspace{3mm}
\noindent\includegraphics[width=\linewidth]{\FIGDIR/f2a_architecture.png}
\figcap{Pooled study architecture over 112 records. \textbf{The largest slice is a defect,
not a design.} \textit{Metadata only} is \textbf{22 records, 20\% of the week} --- surfaced
with no abstract in Crossref or Europe PMC and therefore classified from title and journal
alone --- and it now outranks measurement campaigns, cohorts and every other architecture in
the corpus. It concentrates on two days (six on 19 August, six on 22 August) and on four
Elsevier titles. Against that, \textit{Observational --- cohort} at 15 and
\textit{Experimental / toxicology} at 9 are the load-bearing evidence classes.
\textbf{\textit{Trial / intervention} is 1 record in 112}: a single filtration-intervention
effectiveness study on 19 August. For the fourth week running, no completed randomised
intervention on a PM exposure entered the corpus.}

\clearpage

% -------------------------------------------------------------- pooled analytics
\section*{\sffamily\bfseries\color{Deep}Pooled analytics}
\vspace{-1.5mm}{\color{Teal}\rule{\linewidth}{1.1pt}}\vspace{3mm}

\noindent\includegraphics[width=\linewidth]{\FIGDIR/f1_subtopics.png}
\figcap{Pooled subtopic distribution, 16--22 August. \textit{Exposure assessment \&
modelling} (36) and \textit{Sensing} (22) together are just over half the week;
\textit{Other clinical endpoints} (13) is the largest health cluster, ahead of
\textit{Respiratory \& allergic} (9) and \textit{Burden, policy \& mitigation} (9).
Every one of the ten canonical clusters is occupied, which has now held for four
consecutive weeks.}

\vspace{2mm}
\noindent\includegraphics[width=\linewidth]{\FIGDIR/f2b_endpoint.png}
\figcap{Health endpoint by organ system across the pooled week. \textbf{62 of 112 records
carry no human health endpoint at all} --- 20 because no abstract was deposited and the
endpoint could not be extracted, and the remaining 42 because they are monitoring,
exposure, instrument or aerosol-chemistry work by design. Of what is left, respiratory
endpoints lead at 10, with a further 4 specifically respiratory \emph{infection};
neuro/cognitive and cardiovascular follow.}

\vspace{2mm}
\noindent\includegraphics[width=\linewidth]{\FIGDIR/f3_metric_geography.png}
\figcap{Left --- particle metric across the pooled week. PM$_{2.5}$ alone is the exposure
in 55 records and appears jointly with PM$_{10}$ in 12; \textbf{ultrafine particles are
the explicit subject of 2 records in 112}, and size-resolved number distributions of one
more. That gap has been stable for the entire watch and is the clearest standing mismatch
between what is measured and what deposition physics says matters. Right --- geographic
provenance. \textbf{Europe (25) came within one record of China (26)}, which has not
happened before in this watch; North America is 15. Sub-Saharan Africa and Southeast Asia
contribute two records each.}

\clearpage

\noindent\includegraphics[width=\linewidth]{\FIGDIR/f5_forest.png}

\vspace{1mm}
\noindent\includegraphics[width=\linewidth]{\FIGDIR/f5_forest_2.png}
\figcap{All 41 pooled effect estimates, split across two panels at 32 rows each and
plotted on a common logarithmic axis. \textbf{Exposure contrasts are heterogeneous and the
estimates are not mutually comparable} --- per-IQR, per-10\,$\mu$g/m$^3$, per-SD,
quartile contrasts and binary residential exposures all appear --- so read direction and
precision, never relative magnitude, and \textbf{no meta-analytic summary is computed or
implied}. Two structural features are worth naming. First, \textbf{the week is
front-loaded}: 30 of the 41 estimates come from 16--18 August, 20 August produced none at
all, and 21 August produced one. Second, \textbf{the constituent-resolved studies cluster
tightly just above the null} --- Wu's five childhood-obesity ORs span 1.03--1.06 and Zhao's
seven spontaneous-abortion ORs span 1.14--1.19 --- because a per-unit constituent contrast
is a much smaller exposure step than a quartile or threshold contrast, not because the
effects are weaker. The two estimates left of the null are both from the same paper: the
CHARLS clean-fuel transition.}

\vspace{2mm}
\noindent\includegraphics[width=\linewidth]{\FIGDIR/f6_lifecourse.png}
\figcap{Pooled life-course coverage. Per-stage counts and annotations are summed and
merged across the seven daily issues rather than inherited from any one of them. The
in-utero and early-childhood windows carry the week, consistent with the developmental
mixture cohorts (Wu, Wang, Zhao, Perera) that also supply most of the constituent-resolved
estimates above.}

\vspace{2mm}
\noindent\includegraphics[width=\linewidth]{\FIGDIR/f4_heatmap.png}
\figcap{Subtopic $\times$ architecture over the pooled week. The \textit{Metadata only}
column spreads across most subtopic rows rather than sitting in one corner --- the
abstract-free deposits are not a single topic's problem. \textbf{The structural gap the
watch has reported every week is visible again}: the \textit{Sensing} and \textit{Exposure}
rows contain almost no cells that also carry a health endpoint, and the cohort column
contains almost no cells that measured their own exposure. Instrumentation and
epidemiology remain methodologically separate populations of paper.}

\clearpage

% -------------------------------------------------------------- venues
\section*{\sffamily\bfseries\color{Deep}Venues, provenance and what the pipeline did}
\vspace{-1.5mm}{\color{Teal}\rule{\linewidth}{1.1pt}}\vspace{2.5mm}

\footnotesize
\textbf{Venues.} \textit{Environ Sci Technol} led with 11 records, then
\textit{Atmos Chem Phys} (9), \textit{Atmos Environ} (8) and \textit{Environ Pollut} (6);
\textit{Ecotoxicol Environ Saf}, \textit{Atmos Pollut Res}, \textit{Atmosphere} and
\textit{Atmos Res} contributed 4 each. \textbf{Five preprints} entered the week
(Research Square 3, medRxiv 1, Preprints.org 1) --- 4.5\% of the corpus, and three of the
five carried quantitative results substantial enough to lead a daily issue.

\vspace{1.2mm}
\textbf{Two harvester changes were made inside this period, and both affect what the
figures above show.} On 22 August the supplementary ISSN list trialled on 20 August was
promoted into the standing Crossref sweep --- 15 additions, each ISSN verified against
Crossref's own journal index, with MDPI \textit{Sensors} deliberately excluded after
returning 27 deposits for zero in-scope records. The immediate effect was that the
Crossref leg carried 62\% of the 22 August issue while PubMed returned two records. The
immediate cost is the abstract-free share discussed above. \textbf{The first run with
\textit{Environ Sci Technol} in the list also returned 154 works of which 142 were
mastheads and issue-information stubs} being re-deposited in bulk from volumes 42--52;
those carry DOIs and titles and would have passed every downstream check, so a
title-pattern front-matter filter now drops them at the source.

\vspace{1.2mm}
\textbf{Three figure defects were closed during the period and one is disclosed
uncorrected.} The forest-plot axis now enforces a minimum 1.6$\times$ span in log space,
after a single near-null estimate on 21 August was auto-scaled onto an axis spanning 1.5\%
and drawn two-thirds of the way to the right edge. Two design labels
(\texttt{In vitro exposure}, \texttt{In vitro organ-on-chip}) and one more
(\texttt{Health impact model}) were silently falling into ``Other / mixed'' on every issue
and rollup that contained them, and \texttt{Literature review} was mapped to
\textit{Tool / software}; all four are corrected, and this rollup is the first output
rendered with the corrections. \textbf{Disclosed and not corrected:} the endpoint-label
canonicalisation gap reported in the 15 August rollup persists, so \texttt{f2b}'s
``no health endpoint'' bar is assembled here by hand in the caption (62 of 112) rather
than being read off the figure.

\vspace{1.2mm}
\textbf{OpenAlex has now returned HTTP 429 on twenty consecutive runs} and should be
treated as removed from the pipeline rather than degraded. Europe PMC's
\texttt{CREATION\_DATE} leg continues to work but is increasingly dominated by archival
back-scans --- legacy PMC page images with no title and a publication year decades outside
the window --- which accounted for 14 of 28 unique Europe PMC hits on 21 August and 15 of
23 on 22 August. \textbf{The trial registry has not been re-queried since 8 August};
rollups do not query APIs, and the standing trial watch is therefore stale by two weeks
and is not restated in this issue.

\vspace{3mm}
% -------------------------------------------------------------- gaps
\section*{\sffamily\bfseries\color{Deep}Open gaps}
\vspace{-1.5mm}{\color{Teal}\rule{\linewidth}{1.1pt}}\vspace{2.5mm}

\footnotesize
\begin{itemize}
\item \textbf{Nobody measured the exposure they analysed.} Of the 20 tier-A human health
records this week, the overwhelming majority inherit a modelled exposure surface ---
satellite, land-use regression, or a spatiotemporal reanalysis product --- and none reports
that model's error at the spatial scale of its own analysis. The single clear exception is
Liu et al.\ (22 Aug), whose exposure is a directly observed household fuel transition.
This is now the fourth consecutive week the watch has recorded this and it remains the
largest methodological opening for an instrumentation group.

\item \textbf{Constituent-resolved epidemiology has arrived without
constituent-resolved monitoring.} Six papers this week analyse PM$_{2.5}$ by component,
and all six take those components from chemical transport or satellite-derived products.
No record this week measured a speciated ambient time series and linked it to a health
outcome in the same study.

\item \textbf{Ultrafine particles: 2 records in 112.} Unchanged in direction and
magnitude from every previous week of this watch.

\item \textbf{One intervention study in 112 records, and it was not randomised.} The
evidence base being built is almost entirely observational, while the policy questions the
same papers pose --- which precursor to cut, which fuel to subsidise, which filter to
deploy --- are interventional.

\item \textbf{Twenty per cent of the week is unread.} 22 records have no abstract from any
checked source. Two of them (aerosol pH and preterm birth; ammonia-driven ROS generation
in secondary organic aerosol) sit directly on the week's main thread and are the highest
priority for full-text retrieval before the monthly rollup.
\end{itemize}

\clearpage

% -------------------------------------------------------------- digest
\section*{\sffamily\bfseries\color{Deep}Paper-level digest --- the week by subtopic}
\vspace{-1.5mm}{\color{Teal}\rule{\linewidth}{1.1pt}}\vspace{1mm}

{\fontsize{7.4}{9}\selectfont\color{Slate}Every record that appeared in a daily issue
keeps that issue's critical summary verbatim; nothing here is re-written or re-assessed.
Records surfaced without an abstract are shown as such. Ordered by subtopic, then by the
date of the issue that first carried the record.\par}

\input{digest_body}

\clearpage

% -------------------------------------------------------------- register
\section*{\sffamily\bfseries\color{Deep}Pooled corpus register}
\vspace{-1.5mm}{\color{Teal}\rule{\linewidth}{1.1pt}}\vspace{2.5mm}

{\fontsize{6.7}{8.4}\selectfont
\setlength{\tabcolsep}{2.4pt}
\begin{longtable}{@{}>{\RaggedRight}p{25mm}>{\RaggedRight}p{41mm}>{\RaggedRight}p{27mm}%
>{\RaggedRight}p{23mm}>{\RaggedRight}p{18mm}>{\RaggedRight}p{15mm}@{}}
\toprule
\textsf{\textbf{First author}} & \textsf{\textbf{Focus}} & \textsf{\textbf{Journal}} &
\textsf{\textbf{Design}} & \textsf{\textbf{Metric}} & \textsf{\textbf{Region}}\\
\midrule
\endfirsthead
\toprule
\textsf{\textbf{First author}} & \textsf{\textbf{Focus}} & \textsf{\textbf{Journal}} &
\textsf{\textbf{Design}} & \textsf{\textbf{Metric}} & \textsf{\textbf{Region}}\\
\midrule
\endhead
\multicolumn{6}{@{}l@{}}{\cellcolor{Mist}\textsf{\textbf{\textcolor{Deep}{Neuro / mental health}}}}\\[0.6mm]
\href{https://doi.org/10.1001/jamanetworkopen.2026.25015}{Dales} & Wildfire-sourced PM$_{2.5}$ and ED visits for migraine and other primary headache & \textit{JAMA Netw Open} & Case-crossover & PM$_{2.5}$ (wildfire) & Canada\\
\href{https://doi.org/10.1093/ageing/afag215}{Glover} & Umbrella meta-review of urban environmental exposures and cognitive health & \textit{Age Ageing} & Umbrella review / meta-analysis & PM$_{2.5}$ & Global\\
\href{https://doi.org/10.1016/j.nsa.2026.107019}{Goossens} & In utero air pollution exposure and pre-/postnatal brain development & \textit{Neurosci Appl} & Narrative review & PM$_{2.5}$ / PM$_{10}$ / PM$_{\mathrm{coarse}}$ & Global\\
\href{https://doi.org/10.1016/j.isci.2026.116617}{Ma} & PM$_{2.5}$-induced neurotoxicity and depression: exposure risks and mechanisms & \textit{iScience} & Narrative review & PM$_{2.5}$ + components & Global\\
\href{https://doi.org/10.3760/cma.j.cn112150-20250804-00752}{Pan} & Maternal PM exposure and offspring neurodevelopment: effects and mechanisms & \textit{Zhonghua Yu Fang Yi Xue Za Zhi} & Narrative review & PM (bulk) & China\\
\addlinespace[0.9mm]
\multicolumn{6}{@{}l@{}}{\cellcolor{Mist}\textsf{\textbf{\textcolor{Teal}{Cardiovascular \& metabolic}}}}\\[0.6mm]
\href{https://doi.org/10.3760/cma.j.cn112338-20251120-00843}{Huang} & PM$_{2.5}$ and dyslipidaemia in rural Chinese children; dietary effect modification & \textit{Zhonghua Liu Xing Bing Xue Za Zhi} & Cross-sectional (multistage) & PM$_{2.5}$ & China\\
\href{https://doi.org/10.32394/pe/218966}{Kazmierska} & Air pollution and cardiovascular disease: review of current data & \textit{Przegl Epidemiol} & Narrative review & PM$_{2.5}$ / PM$_{10}$ & Global\\
\href{https://doi.org/10.1007/s00125-026-06814-2}{Xi} & Air pollution, built environment and temperature vs. prediabetes phenotypes & \textit{Diabetologia} & Prospective cohort & PM$_{2.5}$ & Germany\\
\href{https://doi.org/10.1161/JAHA.125.046917}{Yang} & Short-term PM$_{2.5}$ and hospitalisation/ED use in atrial fibrillation & \textit{J Am Heart Assoc} & Case-crossover & PM$_{2.5}$ (wildfire-influenced) & USA\\
\addlinespace[0.9mm]
\multicolumn{6}{@{}l@{}}{\cellcolor{Mist}\textsf{\textbf{\textcolor{Sky}{Reproductive \& developmental}}}}\\[0.6mm]
\href{https://doi.org/10.1016/j.jes.2025.10.043}{Chen} & Prenatal air pollution, fetal growth trajectory phenotypes and child neurodevelopment & \textit{J Environ Sci (China)} & Prospective cohort & PM$_{2.5}$ / PM$_{10}$ & China\\
\href{https://doi.org/10.1016/j.ecoenv.2026.120566}{Huang} & Joint prenatal bisphenol + air-pollutant mixture exposure and preschool IQ & \textit{Ecotoxicol Environ Saf} & Prospective cohort (mixtures) & PM$_{2.5}$ / PM$_{10}$ & China\\
\href{https://doi.org/10.1007/s40572-026-00552-8}{Peterson} & Ambient air pollution, PAHs and pubertal development: critical review & \textit{Curr Environ Health Rep} & Critical review & PM$_{2.5}$ components & Global\\
\href{https://doi.org/10.1016/j.envres.2026.125313}{Xie} & Long-term PM$_{2.5}$ and timing of pubertal development in boys and girls & \textit{Environ Res} & Prospective cohort & PM$_{2.5}$ & China\\
\addlinespace[0.9mm]
\multicolumn{6}{@{}l@{}}{\cellcolor{Mist}\textsf{\textbf{\textcolor{Sage}{Respiratory \& allergic}}}}\\[0.6mm]
\href{https://doi.org/10.1007/s00484-026-03289-z}{Bayar Muluk} & Climate change and otolaryngology: airway disease impacts and emission reduction & \textit{Int J Biometeorol} & Narrative review & PM$_{2.5}$ & Global\\
\href{https://doi.org/10.4103/mgr.MEDGASRES-D-25-00164}{Liao} & Short-term air pollution, respiratory symptoms and lung function in asthma & \textit{Med Gas Res} & Panel study & PM$_{2.5}$ / PM$_{10}$ & China\\
\href{https://doi.org/10.3390/epidemiologia7040097}{Sanchez de Prada} & Ambient PM and asthma-coded admissions in children 0-3 y, Spain & \textit{Epidemiologia} & Ecological panel & PM$_{2.5}$ / PM$_{10}$ & Spain\\
\href{https://doi.org/10.1016/j.jhazmat.2026.143076}{Yang} & Allergenic vs. toxic risk of airborne seed hairs compared with ambient PM$_{2.5}$ & \textit{J Hazard Mater} & Field measurement & PM$_{2.5}$ + bioaerosol & China\\
\addlinespace[0.9mm]
\multicolumn{6}{@{}l@{}}{\cellcolor{Mist}\textsf{\textbf{\textcolor{Violet}{Mechanistic toxicology}}}}\\[0.6mm]
\href{https://doi.org/10.1016/j.jes.2025.12.074}{Ding} & PM$_{2.5}$ promotes vascular calcification via miR-27b-5p/HES1 signalling & \textit{J Environ Sci (China)} & In vitro + organoid + murine & PM$_{2.5}$ & China\\
\href{https://doi.org/10.1007/s00360-026-01698-5}{Maraschi} & Osmoregulatory disruption and metal bioaccumulation from settleable atmospheric PM & \textit{J Comp Physiol B} & Controlled exposure (ecotox) & Settleable APM & Brazil\\
\href{https://doi.org/10.1016/j.jes.2025.11.018}{Wu} & Macrophage NLRP3 inflammasome in real-ambient PM-induced hepatic glucose dysregulation & \textit{J Environ Sci (China)} & Real-ambient murine + co-culture & Real-ambient PM & China\\
\addlinespace[0.9mm]
\multicolumn{6}{@{}l@{}}{\cellcolor{Mist}\textsf{\textbf{\textcolor{Clay}{Occupational \& indoor}}}}\\[0.6mm]
\href{https://doi.org/10.3389/fmedt.2026.1796828}{Chen} & Nested suction instrument to cut intraoperative surgical-smoke PM$_{2.5}$ & \textit{Front Med Technol} & Bench / device experiment & PM$_{2.5}$ (surgical smoke) & China\\
\href{https://doi.org/10.1093/annweh/wxag055}{Cole} & Determinants of respirable crystalline silica exposure in Australian tunnelling & \textit{Ann Work Expo Health} & Observational exposure study & Respirable dust / RCS & Australia\\
\href{https://doi.org/10.1016/j.jhazmat.2026.143081}{Ding} & Stage-resolved hair Zn-isotope dynamics as a firefighting combustion-exposure tracer & \textit{J Hazard Mater} & Repeated-measures biomarker & Combustion-derived PM & China\\
\href{https://doi.org/10.3390/jox16040119}{Ramsden} & Xenobiotic hazards in aircraft cabin air & \textit{J Xenobiot} & Review + risk economics & Ultrafine particles & Global\\
\addlinespace[0.9mm]
\multicolumn{6}{@{}l@{}}{\cellcolor{Mist}\textsf{\textbf{\textcolor{Amber}{Sensing, forecasting \& instrumentation}}}}\\[0.6mm]
\href{https://doi.org/10.1016/j.jenvman.2026.130515}{Shi} & Adaptive spatiotemporal graph transformer for multi-site multi-step PM$_{2.5}$ forecasting & \textit{J Environ Manage} & Deep-learning model & PM$_{2.5}$ & China / India\\
\href{https://doi.org/10.3390/bios16070380}{Tsuchiya} & Seasonal and indoor-outdoor characteristics of fluorescent aerosol particles in offices & \textit{Biosensors} & Real-time sensor campaign & Bioaerosol / coarse & Japan\\
\addlinespace[0.9mm]
\multicolumn{6}{@{}l@{}}{\cellcolor{Mist}\textsf{\textbf{\textcolor{Amber}{Exposure assessment \& modelling}}}}\\[0.6mm]
\href{https://doi.org/10.1029/2025GH001621}{Solorzano-Araque} & Multiscale exposure modelling: LOTOS-EUROS + mobility surveys + particle toxicity & \textit{GeoHealth} & CTM + mobility integration & PM$_{2.5}$ / PM$_{10}$ & Colombia\\
\addlinespace[0.9mm]
\multicolumn{6}{@{}l@{}}{\cellcolor{Mist}\textsf{\textbf{\textcolor{Coral}{Burden, policy \& mitigation}}}}\\[0.6mm]
\href{https://doi.org/10.1002/hsr2.72886}{Karamzad} & Ischaemic heart disease burden in MENA 1990-2021; PM attributable fraction 28.2\% & \textit{Health Sci Rep} & GBD secondary analysis & Ambient PM & MENA\\
\href{https://doi.org/10.1371/journal.pone.0347989}{Rees} & HEPA filtration to reduce respiratory infection in care homes & \textit{PLoS One} & RCT process evaluation & Filtered PM & UK\\
\href{https://doi.org/10.1016/j.jes.2026.01.009}{Shan} & Policy-driven emission reductions and post-pandemic rebound, high-resolution inventories & \textit{J Environ Sci (China)} & Emission inventory & PM emissions & China\\
\addlinespace[0.9mm]
\multicolumn{6}{@{}l@{}}{\cellcolor{Mist}\textsf{\textbf{\textcolor{Slate}{Other clinical endpoints}}}}\\[0.6mm]
\href{https://doi.org/10.3390/dj14070440}{Lupita} & PM$_{2.5}$, preventive-care capacity and caries/edentulism burden across the EU & \textit{Dent J} & Ecological panel & PM$_{2.5}$ & EU-27\\
\href{https://doi.org/10.1016/j.jdin.2026.06.008}{Park} & Long-term PM exposure and risk of androgenetic alopecia & \textit{JAAD Int} & Nationwide cohort & PM$_{2.5}$ / PM$_{10}$ & South Korea\\
\href{https://doi.org/10.3171/2026.3.PEDS2615}{Sun} & Environmental exposure and area deprivation vs. paediatric hydrocephalus outcomes & \textit{J Neurosurg Pediatr} & Registry cohort + geospatial & PM$_{2.5}$ (satellite-derived) & USA\\
\addlinespace[0.9mm]
\bottomrule
\end{longtable}}


\vspace{4mm}
\begin{keybox}
{\sffamily\bfseries\footnotesize\color{Deep}Provenance and method}\\[1.2mm]
{\fontsize{7.2}{8.8}\selectfont
This rollup is an \textbf{aggregation only}. It pools \texttt{state/corpus/*.json} and
\texttt{state/metrics.csv} over \PERIODSTART{}--\PERIODEND{} via
\texttt{PMRW\_START}/\texttt{PMRW\_END} and \textbf{re-queries no literature API}, so every
record was screened and deduplicated in the daily issue that first carried it. Per-record
provenance, source-query windows, exclusions and known metadata defects are documented in
those daily issues and are not restated here.
\textbf{Daily issues covered: 16, 17, 18, 19, 20, 21 and 22 August 2026 --- seven of seven,
no gap in entry-date coverage.} The 21 and 22 August issues were built as a backfill on the
evening of 22 August after an off-schedule 21 August invocation correctly deferred under
the 22:00 rule; the entry windows are contiguous and nothing was skipped.
\textbf{112 records, 112 unique DOIs} --- no duplicate survived the pooling.
\textbf{Known defects carried into this rollup, stated rather than hidden:}
(i) the \texttt{ENDPOINT\_CANON} gap reported on 15 August is unfixed, so \texttt{f2b}
understates its ``no health endpoint'' bar and the caption states the hand-counted figure;
(ii) journal names are free text, so preprint share is a lower bound;
(iii) \textbf{22 records carry no abstract from any checked source and are classified but
unassessed} --- they are counted in every distribution figure and contribute nothing to the
narrative; (iv) the trial registry has not been refreshed since 8 August and no trial
section is included in this issue rather than presenting stale counts as current.
Effect estimates are transcribed from abstracts and use heterogeneous exposure contrasts
--- see the forest-plot caption; \textbf{no meta-analytic summary is computed and none
should be}. Trend and inflection statements are judgements over a seven-day window and are
not statistical tests; where a movement is plausibly an artefact of this pipeline's own
harvesting legs, the text says so. Content derived from publisher metadata; abstracts
remain the copyright of their respective publishers.\par}
\end{keybox}

\end{document}
